\part{Integration guide}\label{um:part:integration-guide}

\chapter{Clock-domain crossing integration}\label{um:ch:cdc-integration}

This chapter gives the tasks for a signal that goes from one clock domain to another. The rules are in \cref{ds:ch:mod-cdc-primitives}. The description of each block is in \cref{tm:ch:mod-cdc-primitives}.

\section{Choose the block}\label{um:ch:cdc-integration:choose}
Find what crosses in \cref{tab:um:ch:cdc-integration:choose}. Use the block of the right column.

\begin{table}[ht]
  \centering
  \caption{Which block to use}\label{tab:um:ch:cdc-integration:choose}
  \begin{tabular}{@{}>{\raggedright\arraybackslash}p{50mm}>{\raggedright\arraybackslash}p{36mm}>{\raggedright\arraybackslash}p{50mm}@{}}
    \toprule
    What crosses & Block & Reason\\
    \midrule
    One bit that stays stable for more than one period of the receiving clock & \code{cdc_sync} & A chain of flops gives the first flop time to settle.\\
    An idle-high line, such as a serial receive pin & \code{cdc_sync} with \code{RESET_VAL} of 1 & The line shows its idle value during reset.\\
    A counter or an index that changes one bit on each source edge & \code{cdc_sync} with \code{WIDTH} above 1 & All bits stay coherent, because only one bit changes.\\
    A multi-bit binary bus & A gray code, or a handshake & A binary bus can show a wrong value when several bits change.\\
    An asynchronous reset for one domain & \code{cdc_reset_sync} & The reset asserts at once and deasserts in step with the clock.\\
    A stream of words with ready and valid & \code{cdc_gray_fifo} & The FIFO keeps the order and gives back-pressure.\\
    The five channels of an AXI4 bus & \code{async_axi_fifo} & It has one FIFO for each channel and a common reset.\\
    Register access by APB4 & \code{apb_cdc_bridge} & It carries one transfer at a time and ends it with an error when the destination is in reset.\\
    A pulse shorter than one receiving period & None of the blocks & Stretch the pulse in the source domain first.\\
    \bottomrule
  \end{tabular}
\end{table}

\section{Instantiate a block}\label{um:ch:cdc-integration:instantiate}
In every task, the names in the listing are examples. Use the names of your design. For \code{cdc_sync} and \code{cdc_gray_fifo}, supply each reset already synchronised to the clock of its domain. The two bridges synchronise their resets inside.

\subsection{A synchroniser for one bit}
\begin{procedure}
  \item Choose \code{STAGES}. The default is 2. A larger value gives a longer MTBF and a longer latency.
  \item Choose \code{RESET_VAL}. It is the value that the logic after \code{q_o} sees during reset.
  \item Connect \code{clk_i} to the clock of the receiving domain, and \code{rst_n_i} to a reset that is synchronised to this clock.
  \item Add the instance to the file list of every flow that uses it (\cref{um:ch:cdc-integration:verify}).
\end{procedure}
\begin{lstlisting}
cdc_sync #(
    .WIDTH     (1),
    .STAGES    (2),
    .RESET_VAL (1'b1)
) u_rx_sync (
    .clk_i   (clk_i),
    .rst_n_i (rst_n_i),
    .d_i     (uart_rx_i),
    .q_o     (rx_sync)
);
\end{lstlisting}

\subsection{A reset synchroniser}
\begin{procedure}
  \item Connect \code{clk_i} to the clock of the domain that gets the reset, and \code{rst_n_i} to the source reset.
  \item Connect \code{scanmode_i} and \code{scan_rst_ni} of the scan flow. Without scan, tie \code{scanmode_i} to \code{1'b0} and \code{scan_rst_ni} to \code{1'b1}.
  \item Hold the source reset low for at least the minimum pulse width of the flops. The chain clears without a clock edge.
\end{procedure}
\begin{lstlisting}
cdc_reset_sync #(
    .STAGES (2)
) u_dom_rst_sync (
    .clk_i       (dom_clk),
    .rst_n_i     (src_rst_n),
    .scanmode_i  (1'b0),
    .scan_rst_ni (1'b1),
    .rst_n_o     (dom_rst_n)
);
\end{lstlisting}

\subsection{A FIFO between two domains}
\begin{procedure}
  \item Choose \code{WIDTH}, \code{DEPTH} (a power of two, at least 2) and \code{SYNC_STAGES}.
  \item Connect the write side to the write clock and a reset that is synchronised to it. Connect the read side in the same way.
  \item Take an entry when \code{rd_valid_o} and \code{rd_ready_i} are both high. Offer an entry with \code{wr_valid_i} and wait for \code{wr_ready_o}.
  \item Keep the instance names \code{u_wr_ptr_to_rd} and \code{u_rd_ptr_to_wr} inside the module, and add a delay budget for the new instance (\cref{um:ch:cdc-integration:constrain}).
\end{procedure}
\begin{lstlisting}
cdc_gray_fifo #(
    .WIDTH       (16),
    .DEPTH       (4),
    .SYNC_STAGES (2)
) u_xing_fifo (
    .wr_clk_i   (a_clk),   .wr_rst_n_i (a_rst_n),
    .wr_data_i  (a_data),  .wr_valid_i (a_valid), .wr_ready_o (a_ready),
    .rd_clk_i   (b_clk),   .rd_rst_n_i (b_rst_n),
    .rd_data_o  (b_data),  .rd_valid_o (b_valid), .rd_ready_i (b_ready)
);
\end{lstlisting}

\subsection{An APB bridge}
\begin{procedure}
  \item Connect the \code{s_*} face to the APB master and the \code{m_*} face to the APB slave. Set \code{ADDR_W} to the address width of both.
  \item Connect \code{s_pstrb_i}. If the master has no strobes, tie it to \code{4'hF}. Leave \code{m_pstrb_o} open if the slave has no strobes.
  \item Connect the scan pins as for the reset synchroniser.
  \item Hold \code{m_rst_n_i} low for more than one period of \code{s_clk_i}. A shorter pulse can make the destination run a transfer again.
  \item Use the bridge for configuration and debug access only. Do not use it for bulk data.
\end{procedure}
\begin{lstlisting}
apb_cdc_bridge #(
    .ADDR_W (12)
) u_cfg_cdc (
    .scanmode_i (1'b0), .scan_rst_ni (1'b1),
    .s_clk_i (core_clk), .s_rst_n_i (core_rst_n),
    .s_psel_i (psel), .s_penable_i (penable), .s_pwrite_i (pwrite),
    .s_paddr_i (paddr), .s_pwdata_i (pwdata), .s_pstrb_i (4'hF),
    .s_prdata_o (prdata), .s_pready_o (pready), .s_pslverr_o (pslverr),
    .m_clk_i (dev_clk), .m_rst_n_i (dev_rst_n),
    .m_psel_o (m_psel), .m_penable_o (m_penable), .m_pwrite_o (m_pwrite),
    .m_paddr_o (m_paddr), .m_pwdata_o (m_pwdata), .m_pstrb_o (m_pstrb),
    .m_prdata_i (m_prdata), .m_pready_i (m_pready), .m_pslverr_i (m_pslverr)
);
\end{lstlisting}

\subsection{An AXI4 bridge}
\begin{procedure}
  \item Connect the \code{s_*} face to the AXI4 master and the \code{m_*} face to the AXI4 slave. The faces have no identifier.
  \item Take \code{s_rst_n_i} and \code{m_rst_n_i} from a common root that also resets the master and the slave side. A reset of one input flushes both faces. In \code{soc_top} the two resets of \code{u_cpu_axi_cdc} come from two top-level ports, so the common root is a task of the board or of the reset controller.
  \item In \code{soc_top}, use \code{cpu_core_rst_n} for \code{s_rst_n_i} of \code{u_cpu_axi_cdc}, and not \code{cpu_domain_rst_n}. The domain reset also drops for a PMU reset, and a flush would lose a response that the CPU waits for.
  \item Check that the clock of each face runs while its reset is low.
  \item Connect the scan pins as for the reset synchroniser.
\end{procedure}
\begin{warningbox}
The bridge \code{apb_cdc_bridge} has these limits. Plan for them.
\begin{itemize}
  \item An error from the bridge during a reset of the destination does not prove that the destination made no transfer. A write can have taken place. Treat such an error as ``access state unknown''.
  \item The bridge has no timeout. When \code{m_clk_i} stops and \code{m_rst_n_i} stays high, the transfer waits, and the APB master upstream waits with it. In \code{soc_top} the \code{m_*} face of \code{u_apb_dbg_cdc} runs on \code{cpu_gated_clk}, which the PMU can stop. Do not stop \code{m_clk_i} while a transfer can be in flight, or assert \code{m_rst_n_i} before you stop it.
  \item A pulse on \code{m_rst_n_i} that is too short can make the destination run a transfer again.
\end{itemize}
\end{warningbox}
\begin{warningbox}
The bridge \code{async_axi_fifo} has these limits. While a response channel (B or R) is empty, its payload outputs on the \code{s_*} face can change on an edge of \code{m_clk_i}. Use the payload only while the valid output is high. A reset of one face also withdraws the valid outputs of the other face. In \code{soc_top}, a CPU transaction in the middle of a burst at clock gating can wedge a crossbar slave engine, because no drain step exists.
\end{warningbox}
The port list of \code{async_axi_fifo} is long and follows the AXI4 names. \Cref{tab:tm:mod:async-axi-fifo:signals} lists the ports.

\section{Constrain the crossings}\label{um:ch:cdc-integration:constrain}
Every path that crosses a clock domain without a synchroniser needs a bounded delay. \Cref{tm:ch:mod-cdc-primitives:constraints} tells why the project uses two constraint files.

\begin{procedure}
  \item Keep the sign-off constraint file as it is. It declares the clocks asynchronous and has no crossing budget.
  \item Open the check constraint file of the flow: \src{pnr/constraints/phase5_soc_multiclock_check.sdc} for ASAP7, or \src{pnr/sky130/soc/constraints/sky130_soc_cdc_check.sdc} for Sky130.
  \item Add a budget for each registered pointer or level that leaves the new instance. Copy the form of the existing part for \code{u_cpu_axi_cdc} (FIFOs) or for \code{u_apb_pll2_cdc} (APB bridge). Use a bound of at most one period of the receiving clock.
  \item Keep the part of the file that comes before the parts for the single instances. The last exception that you declare wins when two exceptions overlap.
  \item Keep the names that the files match: \code{req_q}, \code{ack_q}, \code{u_wr_ptr_to_rd}, \code{u_rd_ptr_to_wr}, \code{rd_gray_q}, \code{sync_q}, and the text \code{rst_sync} in the name of each reset synchroniser instance. The ASAP7 check file matches \code{req_q}, \code{ack_q}, \code{u_wr_ptr_to_rd}, \code{rd_gray_q} and \code{rst_sync}. The Sky130 check file matches \code{u_wr_ptr_to_rd}, \code{u_rd_ptr_to_wr}, \code{u_req_sync}, \code{u_ack_sync} and \code{sync_q}. The CDC configuration matches \code{sync_q} and the instance names.
  \item Run the timing check (\cref{um:ch:cdc-integration:verify}). A line that contains \code{CDC-EXCEPTION-MISS} means that a name pattern matched nothing. Fix the pattern. A budget that is not applied leaves the crossing without a constraint.
\end{procedure}
The bridge \code{u_apb_pll_cdc} has no budget and no CDC check, because the constraint files and the CDC configuration put both of its clocks in the domain of \code{clk_i}. Do not rely on a tool for this crossing. The check file is not a sign-off file. Use it for targeted queries only. A full sweep of the file shows false violations on other pairs of flops.

\section{Verify the crossings}\label{um:ch:cdc-integration:verify}
Run the simulation targets through the Nix shell, as for every flow of the project.

\begin{center}
  \begin{tabular}{@{}p{50mm}p{88mm}@{}}
    \toprule
    Block & Command\\
    \midrule
    \code{cdc_sync} & \code{nix develop --command make -C tb/cocotb/soc cdc_sync}\\
    \code{cdc_reset_sync} & \code{nix develop --command make -C tb/cocotb/soc cdc_reset_sync_scan}\\
    \code{cdc_gray_fifo} & \code{nix develop --command make -C tb/cocotb/soc cdc_gray_fifo}\\
    \code{apb_cdc_bridge} & \code{nix develop --command make -C tb/cocotb/soc apb_cdc_bridge}\\
    \code{async_axi_fifo} & \code{nix develop --command make -C tb/cocotb/soc async_axi_fifo}\\
    CPU and fabric together & \code{nix develop --command make -C tb/cocotb/soc soc_multiclock} and \code{soc_multiclock_reset}\\
    \bottomrule
  \end{tabular}
\end{center}
Each target prints a line with the number of tests, passes and failures. A failure also prints the name of the test. The suites are described in \cref{tm:suite:soc-test-cdc-sync}, \cref{tm:suite:soc-test-cdc-reset-sync-scan}, \cref{tm:suite:soc-test-apb-cdc-bridge} and \cref{tm:suite:soc-test-async-axi-fifo}.

Do these tasks for a new instance.
\begin{procedure}
  \item Add the new module to the source lists of the flows, as the project guide requires for the sv2v file lists. A module that is missing from a list becomes a black box without an error.
  \item Run the CDC check. Inside the \code{sim} directory, run \code{make cdc}. The run is informational by default, and the target returns 0 even when the gate fails. Add \code{CDC_STRICT=1} to make the run fail on an unwaived crossing. To run only the gate on an existing report, run \code{tools/cdc/cdc_gate.py sim/build/cdc/soc_top_cdc.txt -v}. Use the yosys version that \src{tools/cdc/versions.env} pins.
  \item For a new synchroniser that guards a pin, add an entry to the list of required synchronisers in \src{tools/cdc/cdc_config.yml}. Without an entry, a deleted synchroniser stays unseen.
  \item Run the timing check. For ASAP7, run \code{make -C pnr check-cdc-timing-asap7-soc RUN_DIR=<run directory>} on a routed run. For Sky130, run the target \code{check-cdc-timing-sky130-soc}. The make target fails on a finding and on a setup error. Read the lines that start with \code{marker:} to tell them apart. The script itself exits with 0 for a pass, 1 for a finding and 2 for a setup error.
  \item Run the unit tests of the gate: \code{python -m pytest tb/tests/test_cdc_gate.py}.
\end{procedure}
To check a change of \code{cdc_gray_fifo}, run the mutation campaign inside the simulation shell with \code{python3 tools/verif/mutate_cdc_gray_fifo.py}. It prints one line for each mutant. The last line is \code{SURVIVORS: none} when the suite killed every mutant. A survivor means that the suite does not check the changed behaviour. Run \code{tools/verif/equiv_cdc_gray_fifo_depth.sh} to compare the working-tree module with the version at \code{OLD_REF}, which is \code{origin/main} by default, at the depths 4, 8 and 16. On a new host, set \code{YOSYS} and \code{SLANG_PLUGIN}.

\section{Checklist}\label{um:ch:cdc-integration:checklist}
Tick each line before the review.
\begin{itemize}
  \item[$\square$] Each signal that crosses a domain goes through one of the five blocks, or has a waiver in the CDC check configuration.
  \item[$\square$] Each reset of a \code{cdc_sync} or \code{cdc_gray_fifo} is synchronised to the clock of its domain.
  \item[$\square$] \code{RESET_VAL} of each \code{cdc_sync} is safe for the logic after it.
  \item[$\square$] Each multi-bit value that crosses \code{cdc_sync} changes one bit at a time.
  \item[$\square$] \code{scanmode_i} and \code{scan_rst_ni} are connected at every instance.
  \item[$\square$] \code{m_rst_n_i} of an APB bridge is held low for more than one period of \code{s_clk_i}.
  \item[$\square$] The two resets of an AXI4 bridge have a common root with the master and the fabric.
  \item[$\square$] The module is in the source lists of all flows.
  \item[$\square$] The names \code{sync_q}, \code{req_q}, \code{ack_q}, \code{u_wr_ptr_to_rd}, \code{u_rd_ptr_to_wr} and \code{rst_sync} are kept where the files match them.
  \item[$\square$] No clock of an APB bridge stops while a transfer can be in flight.
  \item[$\square$] The check constraint file has a budget for each new crossing, and the timing check reports no \code{CDC-EXCEPTION-MISS}.
  \item[$\square$] The CDC check passes, and the new synchronisers are in the list of required synchronisers where they guard a pin.
  \item[$\square$] The test suite of each block that you changed passes.
\end{itemize}
